% This document was generated by an LLM.

\documentclass{essay}

\title{L'H\^opital's Rule: From Calculus to Real Analysis}
\author{}
\date{}

\begin{document}
\maketitle

\section{An elementary calculus proof}

Suppose that \(f\) and \(g\) are continuous on an open interval containing a
real number \(a\), differentiable there except possibly at \(a\), and satisfy
\(f(a)=g(a)=0\). Assume that \(g'(x)\ne0\) whenever \(x\ne a\) is sufficiently
close to \(a\), and that \(f'(x)/g'(x)\) tends to a finite real number \(L\) as
\(x\to a\). We want to prove that \(f(x)/g(x)\) also tends to \(L\).

The bridge between function values and derivatives is Cauchy's Mean Value
Theorem. Take \(x\ne a\) sufficiently close to \(a\). First, \(g(x)\) cannot be
zero: otherwise \(g(a)=g(x)=0\), and Rolle's theorem would produce a point
strictly between \(a\) and \(x\) where \(g'\) vanishes. Cauchy's Mean Value Theorem
therefore gives a point \(c_{x}\) strictly between \(a\) and \(x\) such that
\[ \frac{f(x)}{g(x)} =\frac{f(x)-f(a)}{g(x)-g(a)} =\frac{f'(c_{x})}{g'(c_{x})}. \]
As \(x\to a\), the intermediate point \(c_{x}\) also approaches \(a\), because it
lies between \(a\) and \(x\). The hypothesis on the derivative ratio then
implies that the right-hand side tends to \(L\). Hence
\(\lim_{x\to a}f(x)/g(x)=L\).

This proof contains the essential idea of the finite \(0/0\) form of
L'H\^opital's rule: a ratio of function increments is exactly a ratio of
derivatives at an intermediate point, and that point is forced toward the
limiting point. The notation \(0/0\) describes the separate limits of the
numerator and denominator; no step of the proof divides zero by zero.

The argument is already valid under the assumptions just stated. A more
explicit real-analysis proof will clarify the limiting argument and also
remove an unnecessary restriction: the original functions need not be
defined or continuous at \(a\).

\section{What the elementary proof assumes}

Continuity at \(a\), together with \(f(a)=g(a)=0\), implies that both functions
tend to zero at \(a\). But the converse requires care. Knowing
\(\lim_{x\to a}f(x)=0\) does not tell us that \(f(a)\) exists, or that it equals
zero if it does exist. A limit at \(a\) concerns values at points near \(a\)
other than \(a\) itself.

For example, define \(f(x)=x\) for \(x\ne0\) and \(f(0)=1729\). Then
\(\lim_{x\to0}f(x)=0\), although \(f\) is not continuous at zero. Alternatively,
we could leave \(f(0)\) undefined, and the same limit would still hold. Since
L'H\^opital's rule concerns a limit of a quotient, its natural hypotheses
should not depend on such an irrelevant endpoint value.

Nevertheless, the elementary proof uses \(a\) as an endpoint of a closed
interval. Cauchy's Mean Value Theorem requires continuity on that closed
interval, including its endpoints. This creates an apparent difficulty:
how can we use the theorem when the original functions may not even be
defined at \(a\)? The answer is to construct continuous extensions, apply
the theorem to those extensions, and transfer the conclusion back to the
original functions. Each part of that transfer must be justified.

\section{Continuous extension is a legitimate reduction}

Suppose \(f\) is defined on a punctured neighborhood of \(a\) and
\(\lim_{x\to a}f(x)=0\). Define a new function by
\[
  F(x)=
  \begin{cases}
    f(x),&x\ne a,\\
    0,&x=a.
  \end{cases}
\]
Then \(F\) is continuous at \(a\). Indeed, for every \(\varepsilon>0\), the
assumed limit supplies a \(\delta>0\) such that
\(0<|x-a|<\delta\) implies \(|f(x)|<\varepsilon\). If \(|x-a|<\delta\) and
\(x\ne a\), this gives \(|F(x)-F(a)|<\varepsilon\); if \(x=a\), the same
inequality holds because its left-hand side is zero. These two cases
establish continuity at \(a\) directly from its definition.

This construction does not assume the conclusion of L'H\^opital's rule.
It uses only the already known limit of \(f\) itself. Nor does it assert that
the original \(f\) was continuous at \(a\). The functions \(f\) and \(F\) may have
different domains or different values at \(a\), and that distinction remains
in force throughout the argument.

There are two reasons why replacing \(f\) by \(F\) is useful here. First,
\(F(x)=f(x)\) at every nearby point other than \(a\), so their punctured limits
agree. More generally, two functions that agree throughout a punctured
neighborhood of \(a\) have exactly the same limit behavior at \(a\). This
follows immediately from the definition of limit, whose hypotheses only
test points with \(0<|x-a|<\delta\).

Second, derivatives away from \(a\) are preserved. For any fixed \(t\ne a\)
in the domain, there is a sufficiently small open interval around \(t\)
that does not contain \(a\). The functions \(F\) and \(f\) agree throughout
that interval, so their difference quotients at \(t\) agree for all
sufficiently small nonzero increments. Consequently \(F'(t)=f'(t)\).
Agreement merely at the single point \(t\) would not establish equality
of derivatives; agreement on a neighborhood does.

Construct \(G\) from \(g\) in the same way, assigning \(G(a)=0\). Wherever the
original quotient is defined near \(a\), we have \(F(x)/G(x)=f(x)/g(x)\).
Likewise, \(F'(t)/G'(t)=f'(t)/g'(t)\) for \(t\ne a\) whenever the denominators
are nonzero. Thus the extensions provide the missing endpoint continuity
while preserving the quantities used in the proof.

Nothing here guarantees differentiability at \(a\), and nothing requires
it. For example, the continuous extension of \(f(x)=|x|\) from
\(\mathbb{R}\setminus\{0\}\) is not differentiable at zero. In the mean value
theorem argument, \(a\) is an endpoint; differentiability is required only
in the open interval between the endpoints.

Calling this a continuous extension across a removable discontinuity is
customary when the endpoint value is absent or incorrect. The extension
is possible precisely because the prescribed finite limit already
exists. Changing one value cannot repair the oscillation of
\(\sin(1/(x-a))\) at \(a\), nor can it produce a finite continuous extension
of \(1/(x-a)\) there.

\section{A precise real-analysis theorem}

We now state the result entirely on a punctured interval. This formulation
avoids making any assumptions about endpoint values and specifies where
all the derivatives and quotients are defined.

\begin{theorem}[L'H\^opital's rule, finite \(0/0\) form]
  Let \(a\in\mathbb{R}\), let \(r>0\), and set
  \(D=(a-r,a+r)\setminus\{a\}\). Let \(f,g:D\to\mathbb{R}\) be differentiable,
  and suppose \(g'(x)\ne0\) for every \(x\in D\). Assume that
  \[ \lim_{x\to a}f(x)=\lim_{x\to a}g(x)=0, \qquad \lim_{x\to a}\frac{f'(x)}{g'(x)}=L\in\mathbb{R}. \]
  Then \(g(x)\ne0\) for every \(x\in D\), and
  \(\lim_{x\to a}f(x)/g(x)=L\).
\end{theorem}

Requiring \(g'\) to be nonzero on this whole punctured interval entails no
loss relative to requiring it only sufficiently near \(a\): in that case,
we first shrink \(r\). Neither continuity of the derivatives nor
differentiability at \(a\) is assumed.

\begin{proof}
  Define \(F\) and \(G\) on \((a-r,a+r)\) by retaining the values of \(f\) and \(g\)
  on \(D\) and setting \(F(a)=G(a)=0\). The assumed limits imply that both
  extensions are continuous at \(a\). On \(D\), they are differentiable, with
  \(F'=f'\) and \(G'=g'\), and are therefore continuous there as well.

  For \(x\in D\), let \(J_{x}=[\min\{a,x\},\max\{a,x\}]\). The functions \(F\)
  and \(G\) are continuous on \(J_{x}\) and differentiable on its interior.
  If \(g(x)=0\), then \(G(x)=G(a)=0\), so Rolle's theorem supplies an interior
  point \(\xi\) with \(G'(\xi)=0\). But \(\xi\in D\) and
  \(G'(\xi)=g'(\xi)\ne0\), a contradiction. Thus \(g(x)\ne0\) for every
  \(x\in D\).

  Let \(\varepsilon>0\) be arbitrary. By the assumed limit of the derivative
  ratio, there exists \(\eta>0\) such that, for every \(t\in D\),
  \[ 0<|t-a|<\eta \quad\Longrightarrow\quad \left|\frac{f'(t)}{g'(t)}-L\right|<\varepsilon. \]
  Put \(\delta=\min\{\eta,r\}\), and fix any real \(x\) satisfying
  \(0<|x-a|<\delta\). In particular, \(x\in D\). Cauchy's Mean Value Theorem,
  applied to \(F\) and \(G\) on \(J_{x}\), gives a point \(c\) in the interior of
  \(J_{x}\) for which
  \[ \bigl(F(x)-F(a)\bigr)G'(c) =\bigl(G(x)-G(a)\bigr)F'(c). \]
  This identity has the same orientation for both \(x>a\) and \(x<a\): reversing
  the endpoints changes the sign of both increments. Substituting the
  definitions of the extensions gives \(f(x)g'(c)=g(x)f'(c)\). The factors
  \(g(x)\) and \(g'(c)\) are nonzero, by the preceding argument and the
  hypothesis, respectively. Division is therefore legitimate and yields
  \(f(x)/g(x)=f'(c)/g'(c)\).

  Since \(c\) lies strictly between \(a\) and \(x\), it satisfies
  \(0<|c-a|<|x-a|<\delta\le\eta\). Applying the derivative-ratio estimate
  at \(t=c\) now gives
  \[ \left|\frac{f(x)}{g(x)}-L\right| =\left|\frac{f'(c)}{g'(c)}-L\right| <\varepsilon. \]
  The choice of \(\delta\) depends on \(\varepsilon\), not on \(x\). Since \(x\)
  was arbitrary subject to \(0<|x-a|<\delta\), this proves the required
  limit by its \(\varepsilon\)--\(\delta\) definition.
\end{proof}

\section{What changed between the proofs?}

The mathematical mechanism did not change. Both proofs use Cauchy's Mean
Value Theorem to represent a quotient of increments by a quotient of
derivatives at one intermediate point. The second proof separates the
steps that make this representation legitimate and explicitly shows how
the limit estimate is transferred.

The first change concerns the statement being proved. The elementary
version assumes continuity at \(a\) and zero endpoint values. The second
version assumes only zero limits and differentiability away from \(a\).
It therefore covers functions whose endpoint values are undefined or
arbitrary. Continuous extension reduces this more general situation to
the setting of the elementary proof. This reduction could equally well
be included in an introductory calculus proof; it is not a technique
reserved for advanced analysis.

The second change concerns division. Cauchy's Mean Value Theorem initially
gives a cross-multiplied identity. The quotient form follows only after
checking the relevant denominators. The condition \(g'(x)\ne0\) supplies one
denominator directly, while Rolle's theorem, combined with the zero
limit of \(g\), supplies the other. The zero-limit condition matters here:
a function can have a nowhere-zero derivative and still vanish somewhere,
as \(g(x)=x-1\) illustrates.

The third change concerns the intermediate point. Writing \(c_{x}\to a\) is
a concise way to express a valid argument, but its justification is the
inequality \(0<|c_{x}-a|<|x-a|\). The point generally depends on \(x\), and
Cauchy's theorem need not determine it uniquely. We neither assume nor
need continuity, measurability, or an explicit formula for a selection
\(x\mapsto c_{x}\).

In fact, the detailed proof does not construct such a function at all.
It fixes an arbitrary \(x\), obtains one suitable \(c\), and applies an
estimate that holds for every point sufficiently close to \(a\). The
quantifier order is essential: first choose \(\delta\) from
\(\varepsilon\), then consider any eligible \(x\), and only then use the
existence of \(c\). The same \(\delta\) works for every such \(x\) because
every possible intermediate point lies inside the required neighborhood.

It would be misleading to identify rigor with educational level or with
the number of displayed quantifiers. This is an elementary real-analysis
theorem with a fully rigorous proof; no specifically graduate machinery
is needed. An experienced mathematician might compress the second proof
into a paragraph, leaving standard deductions implicit. What matters is
that the deductions are justified, not that every proof spells out each
one at maximum length.

\section{Why the common intermediate point matters}

Applying the ordinary Mean Value Theorem separately to \(F\) and \(G\) would
give points \(\xi_{x}\) and \(\zeta_{x}\) between \(a\) and \(x\) such that
\(f(x)=f'(\xi_{x})(x-a)\) and \(g(x)=g'(\zeta_{x})(x-a)\). After division, this
would yield \(f(x)/g(x)=f'(\xi_{x})/g'(\zeta_{x})\). Both intermediate points
approach \(a\), but the hypothesis controls \(f'(t)/g'(t)\) at a
\emph{common} argument. It does not in general control ratios evaluated
at two different arguments.

For a concrete illustration, take \(f(t)=g(t)=t^{2}\) and \(a=0\). The
same-argument derivative ratio is identically one on the punctured
domain. Yet for positive \(x\), the two points \(u_{x}=x/2\) and \(v_{x}=x/3\)
both lie between zero and \(x\), while \(f'(u_{x})/g'(v_{x})=3/2\). These points
are not asserted to be the mean value points for those particular
increments; the example shows exactly why their location and convergence
alone cannot justify the desired inference.

Cauchy's theorem supplies the additional relation that the separate
applications do not: one point serves both derivatives. Its exact
identity is stronger than a heuristic assertion that sufficiently small
increments are approximately derivatives times the same displacement.
The latter heuristic needs further error control, especially when the
denominator derivative approaches zero.

\section{Scope and logical direction}

The theorem proved here concerns a finite limiting point, zero limits
for both functions, and a finite limit of the derivative ratio. The same
endpoint argument gives one-sided versions by working only on the
appropriate side of \(a\). Infinite derivative-ratio limits can be treated
by replacing the \(\varepsilon\) estimate with the defining threshold
inequalities for divergence to infinity. Versions involving unbounded
denominators or a limiting point at infinity require additional arguments;
they are not established merely by the continuous-extension proof above.

L'H\^opital's rule is a sufficient criterion, not an equivalence. Even in
the finite \(0/0\) setting, the quotient may converge while the derivative
ratio does not. For \(x\ne0\), take \(f(x)=x+x^{2}\sin(1/x)\) and \(g(x)=x\).
Both tend to zero and \(g'(x)=1\), while
\[ \frac{f(x)}{g(x)}=1+x\sin(1/x)\longrightarrow1, \qquad \frac{f'(x)}{g'(x)}=1+2x\sin(1/x)-\cos(1/x). \]
The derivative ratio equals zero along \(x_{n}=1/(2\pi n)\) and two along
\(y_{n}=1/((2n+1)\pi)\), for positive integers \(n\), so it has no limit at
zero. Failure of the derivative-ratio hypothesis therefore leaves the
original limit undecided.

The central lesson of the proof is the relationship between local
information and a limit. We know how the derivative ratio behaves at
every sufficiently nearby punctured point. Cauchy's theorem identifies
the original quotient with that ratio at one such point. Continuous
extension makes the theorem applicable at the endpoint without altering
either the original quotient nearby or the derivatives used to estimate
it. Every modification serves a stated hypothesis, and every conclusion
about an auxiliary function is transferred through an exact equality.

\end{document}
